\documentclass[12pt]{report}
\usepackage[a4paper,margin=1in]{geometry}
\usepackage[T1]{fontenc}
\usepackage[utf8]{inputenc}
\usepackage{lmodern}
\usepackage{setspace}
\usepackage{amsmath,amssymb}
\usepackage{booktabs}
\usepackage{multirow}
\usepackage{array}
\usepackage{graphicx}
\usepackage{float}
\usepackage{caption}
\usepackage{subcaption}
\usepackage{hyperref}
\usepackage{enumitem}
\usepackage{longtable}
\usepackage{pdflscape}
\usepackage[table]{xcolor}
\usepackage{url}
\usepackage[strings]{underscore}
\usepackage{fancyhdr}
\usepackage{titlesec}
\usepackage{tabularx}
\usepackage{makecell}
\usepackage{ragged2e}
\usepackage[numbers]{natbib}
\hypersetup{colorlinks=true,linkcolor=blue,citecolor=blue,urlcolor=blue}
\setlength{\parindent}{0pt}
\setlength{\parskip}{0.6em}
\onehalfspacing
\pagestyle{fancy}
\fancyhf{}
\fancyhead[R]{\thepage}
\setlength{\headheight}{15pt}
\fancyhead[L]{Next-Tick IV Curve Forecasting}
\renewcommand{\headrulewidth}{0.4pt}
\titleformat{\chapter}[display]{\bfseries\Huge}{\chaptername\ \thechapter}{0.4em}{\LARGE}
\titleformat{\section}{\bfseries\Large}{\thesection}{0.6em}{}
\titleformat{\subsection}{\bfseries\large}{\thesubsection}{0.6em}{}

\newcommand{\assetroot}{thesis_report_assets}

\newcommand{\insertassetfigure}[3]{%
\begin{figure}[H]
\centering
\IfFileExists{\assetroot/figures/#3}{\includegraphics[width=0.94\textwidth]{\assetroot/figures/#3}}{\fbox{\parbox[c][2.2in][c]{0.88\textwidth}{\centering Missing figure\\[0.4em]\texttt{\assetroot/figures/#3}}}}
\caption{#1}
\label{#2}
\end{figure}}

\begin{document}

\begin{titlepage}
\centering
{\Large National University of Singapore\par}
\vspace{1.5cm}
{\Huge\bfseries Next-Tick Implied-Volatility Curve Forecasting\\[0.3em]for Options Mispricing Detection\par}
\vspace{1cm}
{\Large Final Report\par}
\vspace{2cm}
{\large James Shen\par}
{\large School of Computing\par}
\vfill
{\large March 2026\par}
\end{titlepage}

\pagenumbering{roman}
\chapter*{Abstract}
\addcontentsline{toc}{chapter}{Abstract}
This report studies next-tick implied-volatility (IV) curve forecasting for SPY options on a fixed seven-point moneyness grid and a single maturity bucket centered around 30 days to expiry. The forecasting task is to predict the next IV curve, compare the forecast with the currently quoted curve, and interpret the discrepancy as a stylized options-mispricing signal. The final system combines provider abstraction, fixed-grid IV panel construction, a broad baseline library, recurrent and non-recurrent neural models, walk-forward validation, thresholded signal generation, and execution-aware economic evaluation.

The experimental program progressed from a daily minimum viable pipeline to hourly studies, then to stitched walk-forward evaluation, and finally to a larger 5-minute intraday study with richer features, multiple horizons, capacity sweeps, threshold tuning, standardized common-window comparison, and a cost-aware backtest. Early daily and small-sample hourly settings were dominated by simpler baselines. A tuned LSTM only became compelling once the dataset became larger, the feature set became richer, and the evaluation protocol became stricter. In the 5-minute program, the LSTM first beat matched MLP baselines on standardized common windows at 1-bar, 12-bar, and 24-bar horizons, and these gains were supported by Diebold--Mariano tests, moneyness-bucket analysis, regime splits, execution-stress checks, placebo retraining, and regularization diagnostics. A widened-baseline rerun then showed that the conclusion had to become more precise, and a final overlap-safe two-seed matched benchmark on the \texttt{seq12\_h1} cohort tightened it further: Elastic Net became the best model by mean RMSE, smile-coefficient forecasting became the best model by mean net PnL, xLSTM became the strongest neural model, and the earlier single-seed plain-LSTM headline did not replicate cleanly under rerunning.

The main contribution of the work is therefore not a claim that one recurrent model wins outright. It is an end-to-end experimental pipeline showing that claims about IV-curve forecasting depend critically on data granularity, feature construction, baseline strength, temporal validation, seed sensitivity, and economic realism. Under the final protocol, the evidence supports neural sequence models as credible competitors, but it also shows that strong sparse-linear, coefficient-based, and boosting baselines remain at least as important as the neural models in the final ranking.

\tableofcontents
\clearpage
\pagenumbering{arabic}

\chapter{Introduction}

\section{Motivation}
Most deep-learning work in financial forecasting focuses on stock returns or realized volatility. Those targets are easier to standardize, often closer to stationary after common transformations, and more widely available in benchmark datasets. In contrast, this project studies a more structural object: the implied-volatility curve itself. Implied volatility maps directly to option prices through the Greeks, so forecasting IV rather than raw option prices offers a natural route to mispricing analysis. The central question is whether a sequence model can forecast short-horizon IV-curve changes well enough to produce both statistically and economically meaningful improvements over simpler baselines.

This problem is interesting for three reasons. First, the IV curve contains richer cross-sectional structure than a single return or volatility scalar. Second, the target is close to the market object actually used to quote and compare options. Third, very short-horizon IV forecasting is difficult: the signal is small, the data are noisy, and microstructure frictions can erase apparent gains. A credible empirical study cannot stop at one favorable RMSE number. It must show how conclusions change under stronger baselines, stricter temporal validation, standardized test windows, and more realistic economic assumptions.

\section{Research questions}
The study is organized around four research questions:
\begin{enumerate}[leftmargin=1.4cm]
    \item Can a fixed-grid next-tick IV-curve forecasting pipeline be built reliably from the available options and underlying data infrastructure?
    \item Which model families perform best as the problem moves from daily and hourly frequency to richer 5-minute intraday data?
    \item Do any LSTM gains survive stitched walk-forward validation, standardized common-window comparison, and execution-aware economic evaluation?
    \item Are the final gains statistically credible and robust across moneyness regions, market regimes, and stress scenarios?
\end{enumerate}

\section{Main contributions}
The main contributions are:
\begin{enumerate}[leftmargin=1.4cm]
    \item an end-to-end IV-curve forecasting pipeline with provider abstraction, fixed-grid panel construction, baseline and neural models, and a stylized signal-to-execution layer;
    \item a full experimental chronology showing that simple models dominated early low-sample settings, while the LSTM only became compelling after richer intraday data and stricter evaluation;
    \item final 5-minute evidence that matched LSTM finalists beat matched MLP baselines on standardized common windows at 1-bar, 12-bar, and 24-bar horizons, together with an expanded baseline benchmark and a final overlap-safe multi-seed comparison showing where the neural edge survives and where stronger baselines remain ahead;
    \item final robustness evidence from Diebold--Mariano tests, moneyness-bucket analysis, regime splits, execution-stress checks, placebo retraining, vega-weighted-loss ablations, shape-quality diagnostics, best-model-versus-all-baselines comparison on the matched H1 cohort, a multi-seed matched benchmark, and a modern xLSTM extension on the same final 5-minute setup.
\end{enumerate}

\chapter{Literature and Problem Framing}

\section{Where the existing literature concentrates}
A large share of machine-learning finance studies target stock returns, realized volatility, or scalar volatility indices rather than the tradable implied-volatility curve itself. Even within the IV literature, much of the evidence concentrates on daily or lower-frequency surface dynamics, factor extraction, or coefficient forecasting rather than intraday fixed-maturity smile prediction. Bernales and Guidolin study predictable dynamics in equity-option IV surfaces and show that information from the S\&P~500 volatility surface materially improves equity-surface forecasting \citep{bernales2014}. Kearney et al. show that factor-style term-structure models such as Nelson--Siegel remain strong forecasters for commodity implied-volatility surfaces \citep{kearney2019}. More recent work by Chen et al. forecasts the S\&P~500 IV surface by first fitting a B-spline basis and then predicting the time-varying coefficients with tree-based methods \citep{chen2024bspline}. Borochin and Zhao study the economic value of forecasting equity implied-volatility innovations with machine learning and report incremental value from nonlinear models in delta-hedged option settings \citep{borochin2025}. At the same time, Arratia et al. warn that modern IV-forecasting results can be overstated when leakage, feature overlap, or weak temporal splits are not controlled carefully \citep{arratia2025}. Alongside these peer-reviewed references, the initial project framing was also influenced by more applied machine-learning discussions of option-related prediction settings, early LSTM-style financial forecasting exercises, and hybrid GARCH--LSTM volatility modeling ideas \citep{hull2019framing,sullivan2019lstm,zhao2024garchlstm}. Those framing references did not determine the final empirical standard, but they did motivate the early decision to test recurrent models and volatility-prior hybrids rather than relying only on static regressions. The present project therefore focuses on a narrower but more practical setting: intraday forecasting of a fixed-grid IV curve for trading-oriented mispricing detection.

\section{Why IV curves are a meaningful forecasting target}
Implied-volatility surfaces are economically structured objects rather than arbitrary feature vectors. Factor views of the IV surface, such as level, slope, and curvature, motivate reduced-form latent-dynamics models and PCA-based baselines \citep{cont2002dynamics}. The older IV-surface literature already shows that forecasting performance can improve when one models the smile or surface directly rather than collapsing it to a single scalar \citep{bernales2014,kearney2019}. More recent coefficient-based work reinforces the same point from a machine-learning angle: if the surface is encoded through a structured basis, both forecast accuracy and economic interpretation can improve \citep{chen2024bspline}. At the same time, arbitrage-aware parameterizations such as SVI motivate shape-aware smoothing and regularization when fitting or forecasting smiles \citep{gatheral2014arbitrage}. This makes IV curves appealing as a forecasting target: they are closer to the geometry of options pricing than raw returns, but they also require more discipline because forecast shapes should remain economically plausible.

\section{Model classes relevant to this study}
Persistence is a necessary benchmark because IV is strongly autocorrelated at short horizons. AR-style per-grid models test whether local serial dependence is sufficient. Factor AR/VAR models test whether smile dynamics are largely low-dimensional. HAR-style factor models are included because long-memory structure remains a standard benchmark in volatility forecasting and because the IV-surface literature repeatedly finds latent-factor dynamics informative \citep{cont2002dynamics,kearney2019}. GARCH-type models remain natural volatility references because of clustering and time-varying scale, and hybrid volatility-prior discussions in the broader project framing reinforced the decision to keep such baselines in view even once the main study became a curve-forecasting problem \citep{zhao2024garchlstm}. Elastic Net provides a strong sparse linear benchmark. Tree ensembles, gradient boosting, and coefficient-based smile models are important because recent IV-surface papers show that basis-plus-tree or interaction-rich nonlinear methods can be highly competitive without using recurrent networks \citep{chen2024bspline,borochin2025,arratia2025}. MLPs serve as nonlinear, non-recurrent neural competitors. LSTMs remain attractive because they preserve temporal order without flattening the sequence and because recurrent state updates continue to matter in sequence modeling; the earlier project framing also drew on applied LSTM forecasting references that, while not definitive IV-surface evidence, motivated including a simple recurrent benchmark from the start \citep{sullivan2019lstm,hull2019framing}. The later xLSTM extension is motivated by the same question in a more modern recurrent form: if recurrence matters, can a newer memory architecture improve on a tuned plain LSTM in the same IV-forecasting setting \citep{beck2024xlstm}.

\section{Gap addressed by the study}
The final gap addressed here is not merely ``can an LSTM predict IV?'' The more important gap is whether a recurrent model retains value once the forecasting problem is framed in a way that is both tradable and hard to game. The current literature leaves at least four practical gaps. First, much of the work is still daily or lower-frequency, while an options trader interested in mispricing detection cares about intraday re-estimation of a fixed-maturity smile. Second, many studies forecast scalar IV innovations or coefficient vectors rather than a directly tradable fixed-grid curve. Third, strong baseline libraries are often missing or uneven, even though recent papers show that elastic-net, tree-based, and coefficient-forecasting models can remain highly competitive \citep{chen2024bspline,borochin2025,arratia2025}. Fourth, many studies do not simultaneously enforce expanding-window walk-forward validation, standardized common windows, and execution-aware economic filters. A fifth and newer gap is architectural: modern recurrent extensions such as xLSTM have renewed interest in recurrent sequence modeling, but there is essentially no evidence on whether those gains transfer to short-horizon intraday IV-curve forecasting.

This thesis is designed to close exactly that narrower but more practical gap. It studies short-horizon intraday forecasting of a single fixed-maturity seven-point IV curve, translates the forecast into a stylized trading signal, and evaluates the result under leakage-resistant walk-forward retraining, like-for-like common-window clipping, and execution-aware backtesting. The research question is therefore not whether an LSTM can beat a weak benchmark on a convenient split, but whether intraday IV-curve forecasting retains statistical and economic value once the baseline library is widened and the evaluation protocol becomes realistic.

\chapter{Problem Definition and System Design}

\section{Forecasting target}
The final forecasting target is a seven-point implied-volatility curve defined on the fixed moneyness grid
\[
\mathcal{K}=\{-0.15,-0.10,-0.05,0.00,0.05,0.10,0.15\}.
\]
Only a single maturity bucket is modeled, centered around 30 days to expiry. At each timestamp $t$, the panel contains one curve
\[
\sigma_t = (\sigma_t(k_1),\ldots,\sigma_t(k_7)).
\]
The supervised learning problem is
\[
X \in \mathbb{R}^{B\times L\times d}, \qquad y \in \mathbb{R}^{B\times 7},
\]
where $B$ is batch size, $L$ is sequence length, and $d$ is feature dimension.

The target semantics evolved over the project:
\begin{itemize}[leftmargin=1.2cm]
    \item daily MVP: next-day IV curve,
    \item early hourly studies: next-hour IV curve,
    \item medium-horizon hourly study: $+7$ bars,
    \item final 5-minute program: $+1$, $+12$, and $+24$ bars.
\end{itemize}

If the forecast curve at node $k$ is $\hat{\sigma}_{t+h}(k)$ and the current market curve is $\sigma_t^{\mathrm{mkt}}(k)$, the forecast discrepancy is
\[
\Delta IV_t(k)=\hat{\sigma}_{t+h}(k)-\sigma_t^{\mathrm{mkt}}(k).
\]
Following the project framing, the associated price edge is approximated in stylized form by
\[
\Delta P_t(k) \approx \mathrm{Vega}_t(k)\,\Delta IV_t(k).
\]
This is a research proxy rather than a broker-calibrated pricing model. Its purpose is to translate forecast improvements into economically interpretable signals.

\section{Repository architecture}
The codebase was organized modularly so data providers could change without rewriting the modeling layer. The main modules were:
\begin{itemize}[leftmargin=1.2cm]
    \item data layer: provider interfaces, Yahoo and Alpaca providers, CSV/parquet loaders, panel builders, feature engineering, and split logic;
    \item model layer: persistence, AR(1) per grid, factor AR/VAR, GARCH-style baseline, Elastic Net, ExtraTrees, histogram gradient boosting, HAR factor, smile-coefficient, MLP, LSTM, xLSTM, and curve projection utilities;
    \item training layer: baseline and LSTM training scripts, losses, metrics, and early stopping;
    \item evaluation layer: backtesting, statistical tests, strategy diagnostics, and plotting;
    \item experiment scripts: daily, hourly, walk-forward, 5-minute, threshold sweep, standardized comparison, best-model-versus-all-baselines evaluation, and execution-aware diagnostics.
\end{itemize}

\section{Data providers and practical engineering constraints}
The project used multiple providers because free-data constraints were non-trivial. Yahoo Finance was sufficient for daily and short-window hourly underlying data. Alpaca was used as the default options provider and later also as the underlying stock-bar provider when Yahoo's hourly retention limit became restrictive. CSV and parquet fallbacks were added so the modeling layer could operate independently of any one vendor.

A key engineering issue was that Yahoo hourly data retained only about 730 days. This made full long-window hourly backfills from 2024-03-01 to 2026-03-13 unreliable by March 2026. The solution was to implement an Alpaca stock-bar provider for longer-window hourly and 5-minute studies.

\section{IV panel construction}
The models never consumed raw option chains directly. Instead, the system built a fixed-grid IV panel:
\begin{enumerate}[leftmargin=1.2cm]
    \item collect option bars for the target underlying and nearby expiries,
    \item compute or ingest implied volatility,
    \item map strikes to the fixed moneyness grid,
    \item interpolate when needed,
    \item output one timestamped seven-point IV curve.
\end{enumerate}

The final representation remained intentionally simple: one underlying, one maturity bucket, and one fixed grid. This made it possible to iterate quickly on forecasting models before moving to richer surface modeling.

\section{Why SVI was deferred}
An arbitrage-aware smoother such as SVI is a natural next step for IV work, but it was not part of the main panel builder in the final implementation. This was a deliberate trade-off. The practical priority was to build a fully runnable forecasting pipeline and establish whether the signal was present at all. Fixed-grid interpolation was sufficient for that goal. The absence of SVI should be treated as an explicit limitation and future-work direction rather than a hidden weakness.

\chapter{Feature Engineering, Models, and Evaluation Protocol}

\section{Feature engineering evolution}
The initial feature set was intentionally compact. It contained lagged IV-curve values, ATM IV, underlying return, realized-volatility proxies, and DTE bucket information, giving 11 features per timestamp in the early studies. The early datasets contained 260 daily samples, 543 hourly H1 samples, 2045 year-long hourly samples, and 4030 two-year hourly walk-forward samples.

The final 5-minute study expanded the feature set to 19 features because higher-frequency forecasting benefits more from local state variables. The enriched feature set included IV grid values, ATM IV, ATM IV change, underlying return, absolute return, realized volatility, long-window realized volatility, range percentage, log volume, volume z-score, curve slope, curve curvature, and a scaled DTE bucket. The base 5-minute walk-forward dataset contained 35,260 supervised samples, 19 features, and 7 curve outputs. This richer state representation was one of the main reasons the later 5-minute results became materially stronger. Table~\ref{tab:feature_definitions_main} summarizes the final feature families and shows which ones were present in the earlier daily/hourly studies and which were only introduced in the final 5-minute program.

\begin{table}[H]
\centering
\caption{Feature groups used across the study.}
\label{tab:feature_definitions_main}
\scriptsize
\begin{tabularx}{\textwidth}{p{2.4cm} p{3.4cm} X c c}
\toprule
\textbf{Feature group} & \textbf{Main variables} & \textbf{Why included} & \textbf{Early daily/hourly} & \textbf{Final 5-minute} \\
\midrule
IV grid state & Seven fixed-grid IV nodes & Core cross-sectional smile state; the target is the next curve, so the current curve must be part of the state. & Yes & Yes \\
ATM level and change & ATM IV, ATM IV change & Captures the local level of implied volatility and short-horizon level shocks. & ATM level only & Yes \\
Underlying return state & return, absolute return & Links smile movement to contemporaneous spot shocks and the size of recent price moves. & Partial & Yes \\
Volatility state & realized volatility, long-window realized volatility & Adds short- and medium-horizon volatility regime information. & Short-window only & Yes \\
Trading activity & range percentage, log volume, volume z-score & Captures liquidity and local intraday activity conditions that matter more at 5-minute frequency. & No & Yes \\
Curve shape state & slope, curvature & Encodes level-slope-curvature information directly rather than forcing the model to infer it from nodes alone. & No & Yes \\
Contract state & scaled DTE bucket & Keeps the single-bucket setup explicit and allows later extension to multiple maturity buckets. & Yes & Yes \\
\bottomrule
\end{tabularx}
\end{table}

\section{Baseline models}
\textbf{Persistence.} Predict the next curve as the current curve. This is the essential lower bound because IV is highly autocorrelated.

\textbf{AR(1) per grid point.} Fit a separate autoregressive model at each moneyness node. This captures simple node-wise serial dependence and worked strongly in early low-frequency settings.

\textbf{Factor AR/VAR.} Apply PCA to the curve, forecast the latent factor vector, and reconstruct the curve. This tests whether low-dimensional smile dynamics are sufficient.

\textbf{GARCH-style ATM-shift baseline.} Forecast an ATM-centered volatility shift and propagate it across the curve. This is a volatility-clustering reference, not a full smile model.

\textbf{MLP baseline.} Flatten the recent sequence window and predict the full curve with a feedforward network. This is the main nonlinear non-recurrent competitor.

\textbf{Elastic Net baseline.} Penalized linear regression over the flattened lag window. This became the strongest sparse-linear baseline in the final study and is important because it tests whether most of the short-horizon signal is still linearly recoverable.

\textbf{ExtraTrees baseline.} A nonlinear tree-ensemble benchmark that captures feature interactions without recurrence and remained economically competitive in the final rerun.

\textbf{Histogram Gradient Boosting baseline.} A stronger boosting baseline motivated by recent tree-based IV-surface work. It remained one of the most serious challengers to the LSTM in the final comparisons.

\textbf{HAR factor baseline.} A long-memory factor benchmark motivated by the volatility-forecasting literature and by evidence that low-dimensional IV factors can remain informative.

\textbf{Smile-coefficient baseline.} A structure-preserving baseline that first represents the curve through low-dimensional coefficients and then forecasts those coefficients. This baseline is especially relevant because it stays close to the literature on basis-and-coefficient IV-surface modeling.

\section{LSTM model and structural hooks}
The main neural model is an LSTM encoder over the sequence window with a linear head that outputs the seven-point curve. The implementation included Huber loss, gradient clipping, early stopping, optional smoothness penalties, and optional shape projection via PCA. Placeholders were also built for vega-weighted losses and richer no-arbitrage penalties.

Two regularization findings were particularly important. Removing shape projection could slightly improve RMSE while worsening economic results. Removing the smoothness penalty could slightly improve RMSE while not improving trading performance. This already hinted that point forecast error and trading usefulness are related but not identical objectives.

The final extension added an xLSTM implementation based on the official \texttt{xlstm} package. In this environment the CUDA-only \texttt{sLSTM} path was not available, so the implemented extension used the official \texttt{mLSTM} block stack. This still provided a valid modern recurrent comparison on the same final 5-minute dataset and protocol.

\section{Evaluation protocol}
The strongest methodological contribution of the project was the progression from simple holdout testing to stitched walk-forward evaluation, threshold tuning, standardized common-window comparison, and execution-aware backtesting. The final evaluation logic therefore used four progressively stricter levels:
\begin{enumerate}[leftmargin=1.2cm]
    \item chronological statistical holdout testing,
    \item stitched expanding-window walk-forward evaluation,
    \item threshold tuning for signal activation,
    \item standardized common-window and execution-aware economic evaluation.
\end{enumerate}

The extended backtest supported round-trip costs, position sizing, maximum concurrent positions, gross and net exposure caps, holding-period-aware stacking, and blocked-signal accounting. The default 5-minute execution assumptions used 0.25 per-trade exposure, 0.05 minimum exposure, at most 8 concurrent positions, a gross cap of 2.0, a net cap of 1.0, and 17 bps round-trip cost. The 17 bps cost was decomposed into 0.5 bps fees per side, 4.0 bps half-spread per side, 2.5 bps slippage per side, and 1.5 bps impact per side.

\section{Overfitting controls}
The main anti-overfitting measures were strict chronological ordering of train, validation, and test data, expanding-window walk-forward evaluation, early stopping, train-only feature scaling, structural regularization, standardized common-window comparison, and final significance or placebo checks. These controls do not ``solve'' overfitting in an absolute sense, but they substantially reduce the risk of flattering one lucky configuration.

\chapter{Experimental Program}

\section{Phase 0: offline-safe MVP}
The repository first had to run end-to-end even without live data. This phase validated data assembly, baseline training, LSTM training, plotting, and a toy backtest. It also forced a clean provider abstraction.

\section{Phase 1: daily live SPY MVP}
The first live study used daily SPY underlying data from Yahoo, options from Alpaca, sequence length 10, and a next-day IV-curve target. The dataset contained 260 supervised samples with 11 features.

Results:
\begin{itemize}[leftmargin=1.2cm]
    \item best baseline: AR(1) per grid, RMSE 0.024136,
    \item LSTM: RMSE 0.029209 and $R^2=0.832254$,
    \item backtest net PnL: 0.009021 over 35 trades.
\end{itemize}

This established that the end-to-end path worked, but also showed that the sample was too small for the LSTM to demonstrate a convincing edge.

\section{Phase 2: early hourly experiments}
The next step increased data granularity to hourly. For the immediate next-hour setup, the LSTM obtained RMSE 0.035651 with backtest net PnL 0.021792. This was economically better than the daily setup, but still not enough to claim a forecasting edge over strong baselines.

A medium-horizon $+7$-bar study was then run. An important semantic correction was made: the target was redefined from an end-of-day interpretation to a true rolling fixed-horizon target. Once the backtest was corrected for longer holding periods and overlapping trades, the $+7$-bar strategy performed materially worse than the immediate next-hour setup. The lesson was that a longer horizon did not automatically make the LSTM more useful.

\section{Phase 3: year-long hourly H1 holdout}
The next stronger study used hourly SPY data from 2025-03-01 to 2026-03-13, sequence length 14, target shift 1, and a strict chronological train/validation/test split. The dataset contained 2045 supervised samples.

The best classical baseline was AR(1) per grid with RMSE 0.020711 and $R^2=0.928115$. A two-layer LSTM with hidden size 32 performed better when only the training DataLoader order was shuffled. Under that configuration, the LSTM achieved RMSE 0.020036, $R^2=0.932724$, and net PnL 0.061057.

This was the first stage where the LSTM narrowly beat the strongest classical baseline on a fixed holdout.

\section{Phase 4: hourly capacity sweep}
A capacity sweep was then run over $2\times 32$, $2\times 64$, $2\times 128$, $3\times 32$, $3\times 64$, and $3\times 128$ LSTM architectures.

The best model was the two-layer LSTM with hidden size 128, which achieved RMSE 0.019833, $R^2=0.934081$, and net PnL 0.071932. This suggested that increasing width helped more than increasing depth. The data supported somewhat more capacity, but not arbitrarily deeper recurrence.

\insertassetfigure{Hourly capacity sweep across depth and hidden-size configurations.}{fig:hourly_capacity_sweep}{fig_hourly_capacity_sweep.png}

\section{Phase 5: hourly stitched walk-forward evaluation}
The next major methodological step was stitched walk-forward testing. This study used data from 2024-03-01 to 2026-03-13 with 4030 supervised samples. The walk-forward split used initial train 1227, validation 409, test 409, and step 409, for five stitched folds.

The best hourly baseline at all tested horizons became the MLP rather than AR(1):
\begin{center}
\begin{tabular}{lcc}
\toprule
Horizon & Best baseline RMSE & Best baseline net PnL \\
\midrule
1h & 0.023995 & 0.690198 \\
3h & 0.025890 & 0.813686 \\
7h & 0.027326 & 1.011569 \\
\bottomrule
\end{tabular}
\end{center}

For the LSTM, sequence-length and horizon ablations showed that at 1h with a $2\times 128$ LSTM, sequence length 7 gave the best RMSE, sequence length 14 gave the best PnL, and sequence length 28 was not useful. At fixed sequence length 14, horizon 1h gave the best forecast accuracy, horizon 3h gave the best LSTM trading result, and horizon 7h was too difficult and the baseline won.

This phase changed the interpretation of the project. Under stricter walk-forward testing, the LSTM edge became less obvious and the MLP emerged as a strong competitor. This was a useful anti-overfitting result rather than a failure.

\insertassetfigure{Hourly stitched walk-forward ablation summary, showing how sequence length, forecast horizon, and regularization choices altered the LSTM's statistical and economic performance.}{fig:hourly_walkforward_ablation}{fig_hourly_walkforward_ablation_summary.png}

\section{Phase 6: move to 5-minute data}
The 5-minute study was the largest expansion of the project. It used data from 2024-03-01 to 2026-03-13, Alpaca stock bars at 5-minute frequency, Alpaca options bars at 5-minute frequency, 35,260 supervised samples, 19 features, and walk-forward splits with train 17,630, validation 3,526, test 3,526, and step 3,526.

The LSTM grid covered sequence lengths $\{6,12,48,84,168\}$, horizons $\{1,12,24\}$, and architectures $2\times 128$, $3\times 128$, $2\times 256$, and $3\times 256$. Baselines were evaluated over the same sequence-length and horizon grid.

This was the stage where the LSTM finally showed a broad forecasting edge:
\begin{itemize}[leftmargin=1.2cm]
    \item within the original baseline library, the LSTM beat the best baseline on RMSE in all 15 out of 15 sequence-length and horizon datasets,
    \item the LSTM beat the best baseline on raw net PnL in 8 out of 15 datasets.
\end{itemize}

The best 5-minute LSTM by raw trading PnL was \texttt{seq12\_h1\_l2\_h128}, with RMSE 0.050481 and net PnL 41.900988. The best 5-minute LSTM by RMSE was \texttt{seq84\_h1\_l2\_h256}, with RMSE 0.046091. At this point the real contest was no longer LSTM versus persistence or AR(1), but LSTM versus the strongest baseline in the initial library, which was the MLP.

\insertassetfigure{Full 5-minute grid comparison between the LSTM family and the strongest matched baseline on RMSE. The broad pattern is that the LSTM wins statistically across all tested sequence-length and horizon combinations, even though the best economic configuration is not always the same as the best RMSE configuration.}{fig:five_min_rmse_grid}{fig_5min_lstm_vs_baseline_rmse_grid.png}

\insertassetfigure{5-minute LSTM net PnL grid across sequence lengths and horizons. This figure complements the RMSE grid by showing that the best trading configuration differs from the best pure forecasting configuration.}{fig:five_min_pnl_grid}{fig_5min_pnl_grid.png}

Figures~\ref{fig:five_min_rmse_grid} and~\ref{fig:five_min_pnl_grid} should be read as the broad search stage rather than the final ranking stage. Figure~\ref{fig:five_min_rmse_grid} shows that richer 5-minute data finally gave the LSTM a broad statistical edge in the original model grid, while Figure~\ref{fig:five_min_pnl_grid} shows that the economically best configuration was not automatically the one with the lowest RMSE. This is why the later sections revisit the finalists under stricter, matched comparisons instead of freezing the conclusion at the grid-search stage.

\section{Phase 7: finalists}
Five finalists were retained for further analysis:
\begin{enumerate}[leftmargin=1.2cm]
    \item \texttt{seq12\_h1\_l2\_h128},
    \item \texttt{seq84\_h1\_l2\_h256},
    \item \texttt{seq48\_h12\_l3\_h256},
    \item \texttt{seq168\_h12\_l3\_h256},
    \item \texttt{seq84\_h24\_l3\_h256}.
\end{enumerate}
These finalists covered both forecast-accuracy winners and economic winners across the main horizons of interest.

\section{Phase 8 to 10: execution-aware backtesting, threshold sweep, and standardized comparison}
The original backtest was not strong enough for thesis-level economic interpretation, so the final 5-minute program added execution-aware costs and exposure rules. A threshold sweep was then run over thresholds 0.0005 to 0.0040. The best threshold per finalist increased with horizon: 0.0015 for the strongest 1-bar candidate, around 0.0020--0.0025 for 12-bar candidates, and 0.0030 for the 24-bar finalist. This is consistent with lower signal-to-noise and longer risk carry at longer horizons.

A final fairness problem remained after the large 5-minute sweep: different sequence lengths and horizons implied different stitched test spans. To remove this distortion, a standardized comparison layer intersected prediction timestamps and compared finalists with matched baselines on the same common windows. The first standardized pass was run against matched MLP baselines and established the clean candidate set that later entered the full robustness suite. After the baseline library was widened, the frozen finalists were rerun against stronger refreshed baselines. That second rerun is the final hierarchy that should drive the thesis conclusion.

The refreshed standardized results were:
\begin{center}
\begin{tabular}{lccccc}
\toprule
Finalist & LSTM RMSE & LSTM Net PnL & Best refreshed baseline RMSE & Best refreshed baseline Net PnL & Final reading \\
\midrule
\texttt{seq12\_h1\_l2\_h128} & 0.050481 & 7.198847 & ElasticNet 0.050743 & SmileCoeff 7.129395 & LSTM survives \\
\texttt{seq84\_h1\_l2\_h256} & 0.046091 & 4.518346 & ElasticNet 0.045864 & ElasticNet 4.577378 & baseline wins \\
\texttt{seq48\_h12\_l3\_h256} & 0.046869 & 4.229676 & ElasticNet 0.046649 & ExtraTrees 4.241067 & mixed baseline win \\
\texttt{seq168\_h12\_l3\_h256} & 0.047916 & 4.292157 & ElasticNet 0.047118 & HGB 4.342943 & baseline wins \\
\texttt{seq84\_h24\_l3\_h256} & 0.047626 & 3.078834 & ElasticNet 0.047134 & SmileCoeff 3.128115 & baseline wins \\
\bottomrule
\end{tabular}
\end{center}

These refreshed standardized results changed the hierarchy materially. The earlier MLP-only comparison still matters because it powered the full horizon-wide robustness suite, including the DM, region, regime, placebo, and execution-stress analyses. But once the baseline library was widened, the correct interpretation became narrower: \texttt{seq12\_h1\_l2\_h128} was the strongest single-seed surviving LSTM, not an uncontested final winner. The wider H1 model \texttt{seq84\_h1\_l2\_h256} remained a useful forecasting comparator, but it no longer justified being the headline model once Elastic Net was included.

\insertassetfigure{Threshold-sweep results by horizon. Longer horizons required stricter entry thresholds, which is consistent with weaker signal-to-noise and longer risk carry.}{fig:threshold_sweep}{fig_threshold_sweep_by_horizon.png}

\insertassetfigure{Final overlap-safe standardized benchmark on the matched \texttt{seq12\_h1} cohort. The left panel reports mean RMSE across seeds for the six final models, and the right panel reports mean execution-aware net PnL across the same seeds.}{fig:standardized_comparison}{fig_standardized_common_window_comparison.png}

Figure~\ref{fig:standardized_comparison} is the compact leaderboard view of the final benchmark rather than a single-seed finalist comparison. It shows that the last stage of the project does not end with one dominant winner: Elastic Net is strongest by forecast RMSE, smile-coefficient forecasting is strongest by average net PnL, xLSTM is the strongest neural model, and the rerun plain LSTM is clearly weaker under the strictest protocol. This figure should therefore be read together with Phase~12 as the final benchmark summary rather than as evidence that one model family dominates the others.

\section{Phase 11: xLSTM extension on the matched \texttt{seq12\_h1} setup}
The first architectural extension tested whether a more modern recurrent design could improve on the strongest matched plain LSTM on the same H1 cohort. The xLSTM experiment reused the same \texttt{seq12\_h1} 5-minute dataset, the same walk-forward protocol, and the same standardized comparison layer used for the matched plain-LSTM benchmark. In implementation terms, the extension used the official xLSTM package through its \texttt{mLSTM} block stack with two blocks and embedding dimension 128. The CUDA-only \texttt{sLSTM} path was not available in the current environment, so the reported extension should be interpreted as an official \texttt{mLSTM}-based xLSTM comparison rather than a full mixed xLSTM stack.

The result was informative but not transformative. On the matched \texttt{seq12\_h1} common window:
\begin{center}
\begin{tabular}{lcc}
\toprule
Model & RMSE & Net PnL \\
\midrule
\texttt{seq12\_h1\_l2\_h128} & 0.050481 & 7.198847 \\
\texttt{xlstm\_b2\_e128\_mps} & 0.050968 & 7.133844 \\
\bottomrule
\end{tabular}
\end{center}

At the single-seed stage, the xLSTM remained competitive and still beat weaker baselines such as AR(1), factor AR/VAR, GARCH, persistence, and the original MLP benchmark. However, it did not improve on the tuned plain LSTM, and on the matched \texttt{seq12\_h1} window it also lost to the strongest sparse-linear and tree-based baselines on DM accuracy tests, especially Elastic Net, HAR, and histogram gradient boosting.

\section{Phase 12: overlap-safe multi-seed final benchmark on the matched \texttt{seq12\_h1} cohort}
The strictest final rerun then revisited the matched \texttt{seq12\_h1} setup under multiple seeds. This benchmark compared:
\begin{itemize}[leftmargin=1.2cm]
    \item plain LSTM \texttt{l2\_h128},
    \item xLSTM \texttt{b2\_e128},
    \item Elastic Net,
    \item smile-coefficient,
    \item histogram gradient boosting,
    \item HAR factor.
\end{itemize}

Two completed seeds, 7 and 11, were retained in the final aggregate benchmark. Relative to the earlier multi-seed check, the walk-forward schedule was tightened to use overlapping test windows with half-step advancement: initial train 17,630, validation 3,526, test 3,526, step 1,763, and frontier stitching of the non-overlapping portion of each fold. This produced 7 walk-forward folds while preserving the same final stitched out-of-sample coverage window. The aggregate results were:
\begin{center}
\begin{tabular}{lcc}
\toprule
Model & Mean RMSE & Mean net PnL \\
\midrule
ElasticNet & 0.050735 & 7.049953 \\
HistGradientBoosting & 0.050753 & 6.976482 \\
HARFactor & 0.050756 & 7.058806 \\
xLSTM & 0.050895 & 7.073818 \\
SmileCoefficient & 0.051589 & 7.129831 \\
Plain LSTM & 0.063303 & 3.712089 \\
\bottomrule
\end{tabular}
\end{center}

This overlap-safe final benchmark changed the interpretation of the last stage materially. The earlier single-seed plain-LSTM headline did not replicate cleanly when the same setup was rerun from scratch across seeds. The strongest overall forecasting model in the final rerun remained Elastic Net, the strongest average economic model remained smile-coefficient forecasting, and the strongest neural model remained the xLSTM rather than the plain LSTM. The denser 7-fold protocol slightly improved several models, especially xLSTM and the rerun plain LSTM, but it did not overturn the final hierarchy. The main lesson is not that recurrence failed. It is that the final ranking is sensitive enough that seed-robust rerunning and retraining frequency both matter, and once those are imposed the strongest baselines remain at least as competitive as the neural alternatives.

\chapter{Final Robustness And Diagnostic Evaluation}

\section{Why these robustness tests matter}
The standardized comparison established the clean candidate sets, but a final report should not stop at ``the LSTM has lower RMSE.'' It should also answer whether the edge is statistically significant, where on the curve it comes from, whether it survives different volatility regimes, whether it survives harsher execution assumptions, and whether it persists under placebo and loss-function ablations. The robustness program therefore had two layers. First, a horizon-wide standardized LSTM-versus-MLP layer was used across \texttt{h1}, \texttt{h12}, and \texttt{h24} so that the recurrent family could be stress-tested consistently at every horizon. Second, the strongest surviving plain-LSTM H1 specification on the matched \texttt{seq12\_h1} cohort was compared directly against the widened baseline library on the same common window.

\section{Diebold--Mariano significance tests}
Negative Diebold--Mariano statistics favor the LSTM because the loss differential is defined as LSTM loss minus baseline loss. The horizon-wide matched LSTM-versus-MLP comparisons on standardized common windows gave:
\begin{center}
\begin{tabular}{lcccc}
\toprule
Horizon & Overall DM stat & Overall $p$-value & ATM DM stat & ATM $p$-value \\
\midrule
1-bar & -26.317330 & $\approx 0$ & -16.837815 & $\approx 0$ \\
12-bar & -11.903403 & $\approx 0$ & -7.907783 & $\approx 0$ \\
24-bar & -4.665334 & 0.000003 & -8.699618 & $\approx 0$ \\
\bottomrule
\end{tabular}
\end{center}

These results show that the final 5-minute LSTM advantage over the matched MLP is not merely numerical. On the standardized windows, the matched LSTM beat the matched MLP with statistically defensible forecast improvements at all three tested horizons.

A best-model-versus-all-baselines check was run on the matched \texttt{seq12\_h1} window using the widened baseline library, and on that single-seed window the strongest plain LSTM still retained the lowest RMSE and the highest net PnL. The later multi-seed benchmark did not preserve that ranking. The correct interpretation is therefore methodological rather than promotional: one favorable matched window is informative, but the more decisive result comes from seed-robust rerunning under the same protocol.

\insertassetfigure{Summary of the final robustness suite. The figure condenses the standardized-comparison follow-up into significance, moneyness-region, and regime-level evidence.}{fig:robustness_addendum_main}{fig_final_robustness_addendum_summary.png}

\section{Region-level moneyness analysis}
The region-level breakdown was computed for put OTM, ATM, and call OTM regions. The LSTM beat the matched MLP in all three regions at 1-bar and 12-bar horizons. At 24 bars, the LSTM remained better at ATM and on the call side, while the matched MLP was marginally better on the put wing.

\begin{center}
\begin{tabular}{llc}
\toprule
Horizon & Region & RMSE delta (LSTM $-$ MLP) \\
\midrule
1-bar & ATM & -0.003137 \\
1-bar & call OTM & -0.003755 \\
1-bar & put OTM & -0.001526 \\
12-bar & ATM & -0.001534 \\
12-bar & call OTM & -0.002214 \\
12-bar & put OTM & -0.000953 \\
24-bar & ATM & -0.002402 \\
24-bar & call OTM & -0.002055 \\
24-bar & put OTM & 0.000440 \\
\bottomrule
\end{tabular}
\end{center}

This shows that the final edge is not purely an ATM artifact. The most reliable improvements concentrate around ATM and the call side, while the far put wing remains the hardest region, especially at the longest horizon.

\section{Regime robustness}
Regimes were defined by a median split on entry-time ATM IV within each standardized common window. The LSTM beat the matched MLP in both high-ATM-IV and low-ATM-IV regimes at all three horizons.

\begin{center}
\begin{tabular}{llc}
\toprule
Horizon & Regime & RMSE delta (LSTM $-$ MLP) \\
\midrule
1-bar & high ATM IV & -0.002619 \\
1-bar & low ATM IV & -0.002618 \\
12-bar & high ATM IV & -0.001181 \\
12-bar & low ATM IV & -0.001854 \\
24-bar & high ATM IV & -0.000504 \\
24-bar & low ATM IV & -0.001303 \\
\bottomrule
\end{tabular}
\end{center}

This is important because it reduces the risk that the final edge came from one convenient market state only.

\section{Execution sensitivity}
Latency stress was implemented as a latency proxy through higher slippage and market-impact assumptions, because the saved stitched prediction files did not support a full delayed quote-by-quote replay. Even under these harsher assumptions, the strongest LSTM candidates retained positive net PnL, although profits were compressed.

\begin{center}
\begin{tabular}{llcc}
\toprule
Horizon & Scenario & Net PnL & Compression vs base \\
\midrule
1-bar & base & 4.502921 & 0.0\% \\
1-bar & latency proxy & 3.300521 & 26.70\% \\
1-bar & wide spread & 3.601121 & 20.03\% \\
1-bar & wide spread + latency & 2.398721 & 46.73\% \\
\midrule
12-bar & base & 4.220918 & 0.0\% \\
12-bar & latency proxy & 3.465718 & 17.89\% \\
12-bar & wide spread & 3.654518 & 13.42\% \\
12-bar & wide spread + latency & 2.899318 & 31.31\% \\
\midrule
24-bar & base & 3.071533 & 0.0\% \\
24-bar & latency proxy & 2.583333 & 15.89\% \\
24-bar & wide spread & 2.705383 & 11.92\% \\
24-bar & wide spread + latency & 2.217183 & 27.82\% \\
\bottomrule
\end{tabular}
\end{center}

These results reinforce the intended interpretation of the execution-aware layer: harsher implementation assumptions compress raw PnL, but they do not reverse the relative ranking of the strongest candidates.

\insertassetfigure{Execution-sensitivity summary under wider spreads and latency-proxy stress. Absolute profits compress, but the ranking of the strongest finalists remains intact.}{fig:execution_sensitivity_main}{fig_execution_sensitivity.png}

\section{Placebo retraining}
A target-shuffle placebo retraining was run for one LSTM and one MLP. In both families, the placebo condition materially worsened RMSE and reduced economic quality.

\begin{center}
\begin{tabular}{llcccc}
\toprule
Model & Condition & RMSE & Net PnL & Hit rate & Trades \\
\midrule
LSTM & real reference & 0.050481 & 7.198847 & 0.527652 & 8191 \\
LSTM & shuffled target & 0.061146 & 3.948259 & 0.481802 & 10523 \\
MLP & real reference & 0.052500 & 6.871825 & 0.493682 & 8864 \\
MLP & shuffled target & 0.062542 & 3.916340 & 0.477902 & 10544 \\
\bottomrule
\end{tabular}
\end{center}

This is an important safeguard against accidental leakage or spurious structure. If the final edge had been driven by a hidden artifact, the placebo degradation would have been much weaker.

\insertassetfigure{Placebo target-shuffle comparison for one LSTM and one MLP. Both families deteriorate under the placebo condition, which supports the claim that the real models are learning structured signal rather than benefiting from accidental leakage.}{fig:placebo_comparison_main}{fig_placebo_comparison.png}

\section{Vega-weighted-loss ablation}
A vega-weighted-loss ablation was run for the two final H1 LSTM candidates. In both cases, vega-weighting slightly worsened RMSE and did not improve net PnL.

\begin{center}
\begin{tabular}{llccc}
\toprule
Model & Condition & RMSE & Net PnL & Hit rate \\
\midrule
\texttt{seq12\_h1\_l2\_h128} & unweighted & 0.050481 & 7.198847 & 0.527652 \\
\texttt{seq12\_h1\_l2\_h128} & vega weighted & 0.050938 & 7.166402 & 0.516184 \\
\texttt{seq84\_h1\_l2\_h256} & unweighted & 0.046091 & 4.502921 & 0.498503 \\
\texttt{seq84\_h1\_l2\_h256} & vega weighted & 0.046895 & 4.487720 & 0.528027 \\
\bottomrule
\end{tabular}
\end{center}

The result is useful even though it is negative: in the tested finalists, simple unweighted training remained preferable.

\insertassetfigure{Vega-weighted-loss ablation for the two final H1 LSTM candidates. The figure makes clear that vega-weighting did not improve the tested finalists either statistically or economically.}{fig:vega_weighted_ablation_main}{fig_vega_weighted_ablation.png}

\section{Shape-quality diagnostics}
Shape diagnostics were computed for the final H1 forecast family. The key comparison was between the fully regularized model, shape-only, smoothness-only, and unregularized variants. The unregularized and smoothness-free variants had similar RMSE to the base model, but much worse shape-quality indicators, especially multi-kink frequency and slope-sign-change instability.

\begin{center}
\begin{tabular}{lcccc}
\toprule
Variant & RMSE & Net PnL & Multi-kink fraction & Avg. slope sign changes \\
\midrule
base shape + smooth & 0.046091 & 4.502921 & 0.000378 & 0.966534 \\
shape only & 0.046059 & 4.503172 & 0.000378 & 0.968520 \\
smoothness only & 0.046262 & 4.513046 & 0.168936 & 0.931367 \\
unregularized & 0.046074 & 4.501366 & 0.249480 & 1.086595 \\
\bottomrule
\end{tabular}
\end{center}

This supports a central thesis claim: raw RMSE alone is not enough for curve forecasting. Structural regularization helps maintain economically sensible forecast geometry even when headline forecast errors move only slightly.

\insertassetfigure{Shape-diagnostics comparison across the final H1 regularization variants. The unregularized and partially regularized variants stay close on RMSE but produce noticeably worse curve geometry.}{fig:shape_diagnostics_main}{fig_shape_diagnostics_comparison.png}

\chapter{Results Synthesis and Discussion}

\section{What changed across the project}
The project yields four important empirical patterns.

\textbf{First, simple baselines dominated early.} At daily frequency and in the first hourly experiments, autocorrelation was strong enough that persistence and AR-style models were difficult to beat.

\textbf{Second, stricter evaluation reduced overclaiming.} The move from fixed holdout to stitched walk-forward made the LSTM advantage look much weaker on hourly data. This is not a negative outcome. It is evidence that the research process became more honest.

\textbf{Third, richer intraday data changed the model ranking.} At 5-minute frequency, with 19 features and far more samples, the LSTM established a statistically significant and economically meaningful edge over matched MLP baselines. But once the finalist set was rerun against a wider baseline library, the final picture became more nuanced: the LSTM remained strongest in \texttt{seq12\_h1\_l2\_h128}, while Elastic Net, smile-coefficient, tree-ensemble, HAR-style, and gradient-boosting baselines became serious final competitors elsewhere.

\textbf{Fourth, the final ranking remained mixed even after architectural extension.} The xLSTM extension and later overlap-safe multi-seed benchmark showed that modern recurrence can remain competitive on the same data and protocol, and in the final rerun xLSTM became the strongest neural model. But the best overall models in the strictest rerun were still a sparse-linear baseline by RMSE and a smile-coefficient baseline by net PnL.

\section{Why stronger baselines changed the final picture}
The MLP emerged as the strongest classical baseline under the original stitched walk-forward evaluation because the IV forecasting problem is nonlinear, but much of the useful local state can already be captured by flattening recent lags. In moderate sample regimes, MLP training is also easier than recurrent optimization. That is why the first clean benchmark became LSTM versus MLP rather than LSTM versus AR(1).

However, the refreshed finalist rerun showed that MLP was not the only serious comparator. Elastic Net captured a large amount of structured linear signal across the lagged feature window, smile-coefficient forecasting preserved economically meaningful curve geometry, and tree ensembles extracted nonlinear interactions without recurrence. This widened baseline set is what ultimately forced the final thesis conclusion to become more careful.

\section{Why the LSTM only became compelling later}
The LSTM required more samples, richer intraday features, better horizon selection, stronger evaluation discipline, and slightly more capacity before it became competitive. Once those conditions were met, especially in the 5-minute program, the LSTM delivered a consistent forecast advantage. This is the central technical message of the study. The LSTM did not dominate because it was intrinsically superior. It became useful only after the forecasting problem was posed in a way that matched its strengths.

\insertassetfigure{Representative current, forecast, and realized IV curves for the strongest neural final-benchmark model, xLSTM, on the matched H1 setup. This gives a qualitative view of how the recurrent model adjusts smile shape rather than only reporting aggregate error metrics.}{fig:repr_curve_h1_main}{fig_representative_forecast_vs_realized_curve_h1.png}

\section{Forecasting quality versus trading quality}
A repeated finding across the ablations is that lower RMSE does not automatically imply better trading results. Shape projection and smoothness regularization can slightly worsen RMSE while preserving or improving economically usable forecast geometry. The best immediate-horizon 5-minute model by raw trading PnL is not the same as the cleanest model by standardized forecast quality. For that reason, statistical and economic evaluation should be treated as complementary rather than interchangeable.

\section{Most defensible final claim}
The strongest final claim supported by the experiments is the following: at 5-minute frequency, under stitched walk-forward validation, standardized common-window comparison, execution-aware evaluation, and a final overlap-safe multi-seed rerun, neural sequence models are credible competitors for intraday IV-curve forecasting, but the strongest baselines remain at least as competitive. In the strictest final benchmark on the matched \texttt{seq12\_h1} cohort, Elastic Net is best by mean RMSE, smile-coefficient is best by mean net PnL, and xLSTM is the strongest neural model. The earlier single-seed plain-LSTM headline should therefore be interpreted as an informative intermediate result rather than the final winner.

\insertassetfigure{Execution-aware equity curve for xLSTM, the strongest neural model in the final overlap-safe benchmark. The curve is averaged across the retained benchmark seeds so that the figure reflects the final protocol rather than a single favorable run.}{fig:equity_curve_main}{fig_execution_equity_curve_final_model.png}

Figure~\ref{fig:equity_curve_main} should therefore be read as the neural-model counterpart to the benchmark table rather than as proof that xLSTM wins outright. Its value is comparative: it shows that the best recurrent model remains economically credible under the strictest protocol, even though the overall benchmark winners still differ by objective, with Elastic Net leading on RMSE and smile-coefficient forecasting leading on average net PnL.

\chapter{Limitations and Future Work}

\section{Limitations}
The following limitations should be stated clearly.
\begin{enumerate}[leftmargin=1.2cm]
    \item The options panel was reconstructed from free-data infrastructure and is not equivalent to institutional-quality historical surface data.
    \item Sparse intraday strike coverage remained a real issue in the 5-minute panel. A stricter curve-quality filter was prototyped, but it reduced the usable sample too sharply, so the final comparative study retained the broader permissive panel and treats that trade-off explicitly as a limitation.
    \item The curve builder did not yet use SVI or full no-arbitrage smoothing.
    \item The economic layer was execution-aware but still simplified and not a full order-book simulator.
    \item The strategy should be interpreted as a research signal demonstration rather than a production execution system.
    \item The study was restricted to a single underlying and a single maturity bucket.
    \item The IV-to-price edge proxy was intentionally stylized.
    \item The xLSTM extension used the official \texttt{mLSTM} stack only because the CUDA-only \texttt{sLSTM} path was not available in the current environment.
\end{enumerate}

These limitations do not invalidate the main findings. They define the scope of the current contribution and identify the next set of upgrades.

\section{Future work}
Natural extensions include: adding SVI or another arbitrage-aware surface smoother before fixed-grid extraction, extending the work to multiple maturity buckets, incorporating richer quote-side microstructure variables, testing broader xLSTM sweeps or hybrid recurrent-structural models, calibrating execution assumptions with better transaction-cost estimates, evaluating the same protocol on multiple liquid underlyings beyond SPY, and comparing the final recurrent models against richer coefficient-dynamics baselines built directly on cleaner institutional surface data.

\chapter{Conclusion}
This project began as a proposal to forecast next-tick implied-volatility curves with an LSTM and detect mispricings from forecast deviations. The final research outcome is more nuanced and more valuable than that initial framing.

The study first showed that in low-sample daily and hourly settings, simple baselines were hard to beat. It then showed that a tuned LSTM could narrowly outperform classical models on a year-long hourly holdout. Once the methodology was strengthened with stitched walk-forward testing, that early advantage became less convincing and the MLP emerged as a serious competitor. Only after moving to 5-minute data, expanding the feature set, increasing the sample size, and enforcing stronger evaluation discipline did the neural models become genuinely competitive. In the first standardized 5-minute comparison layer, matched LSTM finalists beat matched MLP baselines at 1-bar, 12-bar, and 24-bar horizons, and the advantage was supported by Diebold--Mariano tests, moneyness-bucket analysis, regime robustness, execution sensitivity, placebo retraining, and shape-quality diagnostics. A wider final baseline benchmark then showed that the thesis claim had to become more precise, and the final overlap-safe two-seed matched benchmark tightened it further: Elastic Net became the best model by mean RMSE, smile-coefficient forecasting became the best model by mean net PnL, and xLSTM became the strongest neural model. The earlier single-seed plain-LSTM result therefore remains useful evidence, but not the final winner under the strictest protocol.

The main conclusion is therefore not simply that an LSTM works. It is that next-tick IV-curve forecasting is highly sensitive to data granularity, feature richness, model capacity, baseline strength, seed robustness, and evaluation protocol, and that any claim of model superiority should only be made after careful temporal validation and realistic economic filtering. Under those conditions, the final evidence supports neural sequence models as credible competitors worth taking forward, but the strongest baselines remain at least as important as the neural models in the final ranking.

\appendix

\chapter{Asset-Based Appendix}

This appendix is designed to compile directly once the folder \texttt{thesis_report_assets/} is uploaded alongside this \texttt{.tex} file in Overleaf. Figures are loaded from \texttt{thesis_report_assets/figures/}. The exported CSV tables are intentionally kept external rather than embedded in this document because the full thesis asset pack contains large machine-generated result tables that are better handled as companion material. The asset manifest and source-to-output mapping are available in \texttt{thesis_report_assets/summaries/README\_asset\_manifest.md} and \texttt{thesis_report_assets/summaries/report\_asset\_map.csv}.

\section{Core thesis figures}

\insertassetfigure{Representative forecast versus realized curve: H12.}{fig:app_repr_h12}{fig_representative_forecast_vs_realized_curve_h12.png}
\insertassetfigure{Representative forecast versus realized curve: H24.}{fig:app_repr_h24}{fig_representative_forecast_vs_realized_curve_h24.png}

\section{External table pack}

The following result tables are part of the thesis asset pack and are intended to be submitted alongside the report rather than embedded inside the PDF:

\begin{itemize}[leftmargin=1.2cm]
    \item baseline walk-forward summary,
    \item LSTM ablation summary,
    \item 5-minute LSTM and baseline grid summaries,
    \item finalist threshold and standardized-comparison summaries,
    \item daily and hourly phase summaries,
    \item hourly capacity, sequence-length, horizon, and regularization summaries,
    \item finalists summary,
    \item full robustness and diagnostics table pack.
\end{itemize}

Keeping these machine-generated tables external improves Overleaf reliability and makes the result archive easier to navigate.

\begin{thebibliography}{99}

\bibitem{beck2024xlstm}
Beck, M., P"oppel, K., Spanring, M., Auer, A., Prudnikova, O., Kopp, M., Klambauer, G., Brandstetter, J., and Hochreiter, S. (2024).
\newblock xLSTM: Extended Long Short-Term Memory.
\newblock arXiv preprint arXiv:2405.04517.

\bibitem{arratia2025}
Arratia, A., El Daou, M., Kagerhuber, J., and Smolyarova, Y. (2025).
\newblock Examining Challenges in Implied Volatility Forecasting: A Critical Review of Data Leakage and Feature Engineering combined with High-Complexity Models.
\newblock \emph{Computational Economics}. https://doi.org/10.1007/s10614-025-11172-z

\bibitem{bernales2014}
Bernales, A. and Guidolin, M. (2014).
\newblock Can we forecast the implied volatility surface dynamics of equity options? Predictability and economic value tests.
\newblock \emph{Journal of Banking \& Finance}, 46, 326--342.

\bibitem{borochin2025}
Borochin, P. and Zhao, Y. (2025).
\newblock The economic value of equity implied volatility forecasting with machine learning.
\newblock \emph{Journal of Empirical Finance}, 82.

\bibitem{chen2024bspline}
Chen, Z., Li, Y., and Yu, C.~L. (2024).
\newblock Modeling Implied Volatility Surface Using B-Splines with Time-Dependent Coefficients Predicted by Tree-Based Machine Learning Methods.
\newblock \emph{Mathematics}, 12(7), 1100.

\bibitem{cont2002dynamics}
Cont, R. and da Fonseca, J. (2002).
\newblock Dynamics of implied volatility surfaces.
\newblock \emph{Quantitative Finance}, 2(1), 45--60.

\bibitem{diebold1995}
Diebold, F.~X. and Mariano, R.~S. (1995).
\newblock Comparing predictive accuracy.
\newblock \emph{Journal of Business \& Economic Statistics}, 13(3), 253--263.

\bibitem{gatheral2014arbitrage}
Gatheral, J. and Jacquier, A. (2014).
\newblock Arbitrage-free SVI volatility surfaces.
\newblock \emph{Quantitative Finance}, 14(1), 59--71.

\bibitem{kearney2019}
Kearney, F., Shang, H.~L., and Sheenan, L. (2019).
\newblock Implied volatility surface predictability: The case of commodity markets.
\newblock \emph{Journal of Banking \& Finance}, 108, 105657.

\bibitem{hull2019framing}
Hull, J., Qiao, Y., and others (2019).
\newblock Machine learning in business: implied volatility and option-related prediction settings.
\newblock Referenced as part of the project literature framing.

\bibitem{sullivan2019lstm}
Sullivan, R. (2019).
\newblock LSTM-based financial forecasting course-project style literature referenced in the project framing.
\newblock Referenced in the CA-stage literature review.

\bibitem{zhao2024garchlstm}
Zhao, X., Wang, Y., and coauthors (2024).
\newblock GARCH-LSTM style volatility modeling and neural volatility priors.
\newblock Referenced in the project literature framing.

\end{thebibliography}

\end{document}
